\chapter{USE CASE MODELLING}
\label{ch:use_case_modelling}

\section{Use Case Diagram}
\label{sec:use_case_diagram}

The high-level interactions between the system actors and the core functional use cases of the UniExpert platform are organized into five architectural subsystems in Figure~\ref{fig:use_case_diagram}.

\begin{figure}[H]
\centering
\begin{tikzpicture}[
    scale=0.71, every node/.style={transform shape},
    font=\sffamily\scriptsize,
    actor/.style={rectangle, draw=black!80, thick, fill=blue!6, align=center, minimum width=2.4cm, minimum height=0.95cm, rounded corners=3pt},
    package/.style={rectangle, draw=black!40, dashed, fill=gray!2, rounded corners=4pt, inner sep=5pt},
    pkgtitle/.style={font=\sffamily\bfseries\scriptsize, text=black!70, anchor=north west},
    uc/.style={ellipse, draw=black!80, thick, fill=white, align=center, minimum width=3.3cm, minimum height=0.70cm, font=\sffamily\tiny},
    link/.style={thick, draw=black!70},
    inc/.style={dashed, -{Stealth[scale=0.8]}, draw=black!60}
]

% System Boundary
\draw[thick, rounded corners=8pt, fill=blue!1] (-3.6, -20.7) rectangle (7.6, 1.2);
\node[font=\sffamily\bfseries\small] at (2.0, 0.8) {UniExpert Platform System Boundary};

% Package 1: Identity & Verification (y: 0.3 to -3.2)
\draw[package] (-3.3, -3.2) rectangle (7.3, 0.3);
\node[pkgtitle] at (-3.2, 0.1) {Subsystem 1: Identity \& Verification};
\node[uc] (uc01) at (-1.3, -1.2) {UC01: Register \& Verify Domain};
\node[uc] (uc02) at (4.6, -1.2) {UC02: Log In \& Session Auth};
\node[uc] (uc05) at (4.6, -2.4) {UC05: Verify Academic Tier};

% Package 2: Ingestion & AI Matching (y: -3.5 to -8.0)
\draw[package] (-3.3, -8.0) rectangle (7.3, -3.5);
\node[pkgtitle] at (-3.2, -3.7) {Subsystem 2: NLP Problem Ingestion \& Hybrid Matching};
\node[uc] (uc06) at (-1.3, -4.7) {UC06: Ingest Problem Statement};
\node[uc] (uc07) at (4.6, -4.7) {UC07: Extract Domain \& Skills};
\node[uc] (uc08) at (-1.3, -6.0) {UC08: Refine Skill Weights};
\node[uc] (uc09) at (4.6, -6.0) {UC09: Dense Vector Search};
\node[uc] (uc10) at (-1.3, -7.2) {UC10: Hybrid Match Ranking};
\node[uc] (uc11) at (4.6, -7.2) {UC11: Generate XAI Breakdown};

% Package 3: Profiles & Team Gap Optimization (y: -8.3 to -11.9)
\draw[package] (-3.3, -11.9) rectangle (7.3, -8.3);
\node[pkgtitle] at (-3.2, -8.5) {Subsystem 3: Profile Management \& Team Gap Optimization};
\node[uc] (uc03) at (-1.3, -9.4) {UC03: Manage Student Profile};
\node[uc] (uc04) at (-1.3, -10.8) {UC04: Ingest Expert Profile};
\node[uc] (uc13) at (4.6, -9.4) {UC13: Candidate Comparison Matrix};
\node[uc] (uc14) at (4.6, -10.8) {UC14: Multidisciplinary Gap ($\Delta\mathcal{K}$)};

% Package 4: Collaboration FSM (y: -12.2 to -16.0)
\draw[package] (-3.3, -16.0) rectangle (7.3, -12.2);
\node[pkgtitle] at (-3.2, -12.4) {Subsystem 4: Advisory FSM \& Consultation Workspace};
\node[uc] (uc16) at (-1.3, -13.3) {UC16: Submit Advisory Request};
\node[uc] (uc17) at (4.6, -13.3) {UC17: Advisory FSM Lifecycle};
\node[uc] (uc18) at (-1.3, -14.8) {UC18: Enforce Advisory Quotas};
\node[uc] (uc19) at (4.6, -14.8) {UC19: Schedule Consultation};

% Package 5: Governance & Audit (y: -16.3 to -20.4)
\draw[package] (-3.3, -20.4) rectangle (7.3, -16.3);
\node[pkgtitle] at (-3.2, -16.5) {Subsystem 5: Governance, Audit \& Analytics Reporting};
\node[uc] (uc21) at (-1.3, -17.6) {UC21: Post-Collab Endorsement};
\node[uc] (uc20) at (-1.3, -19.3) {UC20: Workspace Interaction};
\node[uc] (uc22) at (4.6, -17.6) {UC22: Moderate Taxonomy Catalog};
\node[uc] (uc24) at (4.6, -18.7) {UC24: Audit Trail Search \& Export};
\node[uc] (uc25) at (4.6, -19.8) {UC25: Generate Accreditation Report};

% External Actors
% Left Actors:
\node[actor] (student) at (-6.3, -5.3)  {\textbf{Student}\\Project Lead};
\node[actor] (faculty) at (-6.3, -12.8) {\textbf{Faculty}\\Advisor};
\node[actor] (mentor)  at (-6.3, -18.4) {\textbf{Industry}\\Mentor};

% Right Actors:
\node[actor] (admin)   at (9.9, -1.8)   {\textbf{System}\\Administrator};
\node[actor] (coord)   at (9.9, -10.1)  {\textbf{Department}\\Coordinator};
\node[actor] (auditor) at (9.9, -19.2)  {\textbf{Institutional}\\Auditor};

% Direct Connections: Left Actors to Left Column Use Cases (Completely Horizontal / Straight)
\draw[link] (student.east) -- (uc01.west);
\draw[link] (student.east) -- (uc06.west);
\draw[link] (student.east) -- (uc08.west);
\draw[link] (student.east) -- (uc03.west);

\draw[link] (faculty.east) -- (uc04.west);
\draw[link] (faculty.east) -- (uc16.west);
\draw[link] (faculty.east) -- (uc18.west);

\draw[link] (mentor.east) -- (uc21.west);
\draw[link] (mentor.east) -- (uc20.west);

% Direct Connections: Right Actors to Right Column Use Cases (Completely Horizontal / Straight)
\draw[link] (admin.west)   -- (uc02.east);
\draw[link] (admin.west)   -- (uc05.east);

\draw[link] (coord.west)   -- (uc13.east);
\draw[link] (coord.west)   -- (uc14.east);

\draw[link] (auditor.west) -- (uc24.east);
\draw[link] (auditor.west) -- (uc25.east);

% Internal Horizontal Includes (strictly within their subsystem row)
\draw[inc] (uc06) -- node[above, font=\tiny] {$\ll$include$\gg$} (uc07);
\draw[inc] (uc08) -- node[above, font=\tiny] {$\ll$include$\gg$} (uc09);
\draw[inc] (uc10) -- node[above, font=\tiny] {$\ll$include$\gg$} (uc11);
\draw[inc] (uc16) -- node[above, font=\tiny] {$\ll$include$\gg$} (uc17);
\draw[inc] (uc16) -- node[above, font=\tiny, sloped] {$\ll$include$\gg$} (uc19);

\end{tikzpicture}
\caption{System Use Case Diagram for UniExpert Platform}
\label{fig:use_case_diagram}
\end{figure}
\clearpage

\section{Use Case Description}
\label{sec:use_case_descriptions}

The 25 use case descriptions below specify formal behavioral contracts including goals, preconditions, success and failure end conditions, actor roles, triggers, step-by-step main success scenarios, alternative branches, and non-functional quality requirements.

\subsection{Use Case Description-01: Register Account \& Verify Academic Domain}
\label{subsec:uc01}

\begin{table}[H]
\centering
\footnotesize
\begin{tabularx}{\textwidth}{|l|X|}
\hline
\rowcolor{lightgray!40} \textbf{Use Case} & \textbf{UC01: Register Account \& Verify Academic Domain} \\ \hline
\textbf{Goal} & A new student, faculty expert, or industry mentor creates an authenticated account with verified institutional credentials. \\ \hline
\textbf{Preconditions} & The user possesses a valid university email address (\texttt{.edu} or affiliated domain); the registration service is online. \\ \hline
\textbf{Success End Condition} & Account record saved in database; cryptographic verification email dispatched; status \texttt{PENDING\_VERIFICATION}. Message: \textit{``Account registered successfully. Please verify your institutional email.''} \\ \hline
\textbf{Failed End Condition} & No user record created. Message: \textit{``Registration failed: Email address already registered or invalid institutional domain.''} \\ \hline
\textbf{Primary Actors:} & Student Project Lead, Faculty Advisor, Industry Mentor \\
\textbf{Secondary Actors:} & System Administrator \\ \hline
\textbf{Trigger} & User accesses the web portal and clicks ``Register Institutional Account''. \\ \hline
\textbf{Description / Main Success Scenario} &
\begin{enumerate}[leftmargin=*, nosep]
    \item User opens the registration portal and selects their designated institutional role.
    \item User enters full legal name, institutional email, department affiliation, and strong password.
    \item System validates email format against approved university domain whitelist.
    \item System verifies that the email address is not already registered.
    \item System computes Argon2id salted password hash and creates a user record with status \texttt{PENDING\_VERIFICATION}.
    \item System generates a cryptographically signed HMAC token (24-hour TTL) and dispatches an activation link.
    \item User clicks the activation link in their university inbox.
    \item System verifies token authenticity, updates user status to \texttt{ACTIVE}, and redirects to role onboarding.
\end{enumerate} \\ \hline
\textbf{Alternative Flows} &
\begin{itemize}[leftmargin=*, nosep]
    \item[\textbf{3.a.}] \textbf{Email domain is not in university whitelist:} System prompts for valid institutional domain.
    \item[\textbf{4.a.}] \textbf{Email already exists in system:} System displays duplicate notice and provides reset link.
    \item[\textbf{7.a.}] \textbf{Verification token expired (>24 hours):} System offers to resend fresh activation token.
\end{itemize} \\ \hline
\textbf{Quality Requirements} & 
1. Store passwords using salted Argon2id hashes (memory: 64MB, iterations: 3, parallelism: 4). \newline
2. Complete end-to-end registration validation and token generation within 500 milliseconds. \\ \hline
\end{tabularx}
\end{table}

\clearpage
\subsection{Use Case Description-02: Log In \& Session Authentication}
\label{subsec:uc02}

\begin{table}[H]
\centering
\small
\begin{tabularx}{\textwidth}{|l|X|}
\hline
\rowcolor{lightgray!40} \textbf{Use Case} & \textbf{UC02: Log In \& Session Authentication} \\ \hline
\textbf{Goal} & An approved, verified user authenticates credentials and establishes a secure, role-scoped session. \\ \hline
\textbf{Preconditions} & User account is verified and active; account is not currently locked out due to excessive failed attempts. \\ \hline
\textbf{Success End Condition} & User session established; short-lived JWT access token issued; secure HTTP-only refresh token set. Message: \textit{``Authentication successful.''} \\ \hline
\textbf{Failed End Condition} & Access denied; session rejected. Message: \textit{``Invalid email or password. Please try again.''} \\ \hline
\textbf{Primary Actors:} & All Users (Student, Faculty, Mentor, Coordinator, Admin, Auditor) \\
\textbf{Secondary Actors:} & None \\ \hline
\textbf{Trigger} & User navigates to the login screen and submits credentials. \\ \hline
\textbf{Description / Main Success Scenario} &
\begin{enumerate}[leftmargin=*]
    \item User inputs registered email address and password.
    \item System verifies that the account is active and not locked.
    \item System verifies submitted password against stored Argon2id hash.
    \item System checks if multi-factor authentication (MFA) is mandated for user role.
    \item System generates a signed JWT access token containing user UUID, role scopes, and department ID.
    \item System issues an HTTP-only, secure, SameSite refresh token cookie (7-day TTL).
    \item System logs successful authentication event with client IP and User-Agent in the audit trail.
    \item System redirects user to their personalized role dashboard.
\end{enumerate} \\ \hline
\textbf{Alternative Flows} &
\begin{itemize}[leftmargin=*]
    \item[\textbf{2.a.}] \textbf{Account locked due to 5 consecutive failed attempts:}
    \begin{itemize}[leftmargin=*]
        \item[2.a.1.] System rejects login attempt and enforces a 15-minute lockout timer.
    \end{itemize}
    \item[\textbf{3.a.}] \textbf{Password hash mismatch:}
    \begin{itemize}[leftmargin=*]
        \item[3.a.1.] System increments failed attempt counter in Redis and displays invalid credentials error.
    \end{itemize}
    \item[\textbf{4.a.}] \textbf{MFA challenge required (Admin / Coordinator):}
    \begin{itemize}[leftmargin=*]
        \item[4.a.1.] System prompts for 6-digit TOTP code; user submits code; system validates before issuing tokens.
    \end{itemize}
\end{itemize} \\ \hline
\textbf{Quality Requirements} & 
1. Access token expiration enforced at 15 minutes; refresh tokens stored in HttpOnly cookies. \newline
2. Enforce 15-minute account lockout after 5 consecutive failed login attempts. \\ \hline
\end{tabularx}
\end{table}

\clearpage
\subsection{Use Case Description-03: Manage Student Research Profile \& Portfolio}
\label{subsec:uc03}

\begin{table}[H]
\centering
\small
\begin{tabularx}{\textwidth}{|l|X|}
\hline
\rowcolor{lightgray!40} \textbf{Use Case} & \textbf{UC03: Manage Student Research Profile \& Portfolio} \\ \hline
\textbf{Goal} & A student project lead updates team profile details, academic records, CV attachments, and GitHub repositories. \\ \hline
\textbf{Preconditions} & User is signed in with authenticated \texttt{STUDENT} role scopes. \\ \hline
\textbf{Success End Condition} & Profile metadata updated in PostgreSQL; searchable profile card refreshed. Message: \textit{``Student profile updated successfully.''} \\ \hline
\textbf{Failed End Condition} & Profile update rejected. Message: \textit{``Unable to update profile: Invalid input format or file size exceeded.''} \\ \hline
\textbf{Primary Actors:} & Student Project Lead \\
\textbf{Secondary Actors:} & None \\ \hline
\textbf{Trigger} & Student selects ``Edit Profile'' on their dashboard. \\ \hline
\textbf{Description / Main Success Scenario} &
\begin{enumerate}[leftmargin=*]
    \item Student opens the profile management dashboard.
    \item Student inputs degree program, batch number, student ID, bio, and technical competencies.
    \item Student links external portfolio artifacts (GitHub repository URL, LinkedIn profile).
    \item Student uploads an optional PDF resume/CV (max 5 MB).
    \item System validates file integrity, scans for malicious MIME types, and stores PDF in secure object storage.
    \item System persists updated profile entity to PostgreSQL.
    \item System updates profile completion index and notifies student of updated profile score.
\end{enumerate} \\ \hline
\textbf{Alternative Flows} &
\begin{itemize}[leftmargin=*]
    \item[\textbf{4.a.}] \textbf{Uploaded CV exceeds 5 MB or is non-PDF:}
    \begin{itemize}[leftmargin=*]
        \item[4.a.1.] System rejects file upload and prompts student to attach a valid PDF under 5 MB.
    \end{itemize}
    \item[\textbf{6.a.}] \textbf{Database connection timeout during save:}
    \begin{itemize}[leftmargin=*]
        \item[6.a.1.] System rolls back transaction, alerts student, and preserves draft entries in local storage.
    \end{itemize}
\end{itemize} \\ \hline
\textbf{Quality Requirements} & 
1. Restrict CV uploads to authenticated students with maximum file size of 5 MB. \newline
2. Update and re-render profile card in under 1 second. \\ \hline
\end{tabularx}
\end{table}

\clearpage
\subsection{Use Case Description-04: Ingest and Manage Faculty Expert Profile}
\label{subsec:uc04}

\begin{table}[H]
\centering
\small
\begin{tabularx}{\textwidth}{|l|X|}
\hline
\rowcolor{lightgray!40} \textbf{Use Case} & \textbf{UC04: Ingest and Manage Faculty Expert Profile} \\ \hline
\textbf{Goal} & A faculty advisor establishes or synchronizes their research portfolio, publication corpus, and research domains. \\ \hline
\textbf{Preconditions} & User is signed in with \texttt{EXPERT} or \texttt{MENTOR} role scopes. \\ \hline
\textbf{Success End Condition} & Profile synced; publications parsed and indexed; embeddings queued. Message: \textit{``Expert profile and publications updated.''} \\ \hline
\textbf{Failed End Condition} & Synchronization incomplete. Message: \textit{``Publication synchronization failed: External API unreachable.''} \\ \hline
\textbf{Primary Actors:} & Faculty Advisor, Industry Mentor \\
\textbf{Secondary Actors:} & System Administrator \\ \hline
\textbf{Trigger} & Faculty advisor clicks ``Sync Google Scholar / ORCID'' or manually edits research interests. \\ \hline
\textbf{Description / Main Success Scenario} &
\begin{enumerate}[leftmargin=*]
    \item Faculty advisor navigates to the Research Portfolio settings.
    \item Advisor enters their public ORCID iD or Google Scholar profile identifier.
    \item System initiates asynchronous background worker to fetch publication metadata, abstracts, and citations.
    \item Background worker parses ingested papers, computes h-index, and filters top research topics.
    \item System compiles a consolidated academic bio-corpus for the advisor.
    \item System updates advisor record in PostgreSQL and triggers vector re-indexing worker.
    \item System notifies faculty advisor upon completion of publication ingestion.
\end{enumerate} \\ \hline
\textbf{Alternative Flows} &
\begin{itemize}[leftmargin=*]
    \item[\textbf{3.a.}] \textbf{External Scholar API rate limit or network timeout:}
    \begin{itemize}[leftmargin=*]
        \item[3.a.1.] System retries request with exponential backoff (up to 3 retries).
        \item[3.a.2.] If unsuccessful, system allows manual CSV publication upload.
    \end{itemize}
    \item[\textbf{4.a.}] \textbf{No publications detected for entered profile ID:}
    \begin{itemize}[leftmargin=*]
        \item[4.a.1.] System prompts advisor to verify ID or manually input published paper titles.
    \end{itemize}
\end{itemize} \\ \hline
\textbf{Quality Requirements} & 
1. Ingest up to 50 publication abstracts within 10 seconds of background job trigger. \newline
2. Preserve existing verified publications if external sync job fails. \\ \hline
\end{tabularx}
\end{table}

\clearpage
\subsection{Use Case Description-05: Verify Faculty Academic Credentials \& Tier Assignment}
\label{subsec:uc05}

\begin{table}[H]
\centering
\small
\begin{tabularx}{\textwidth}{|l|X|}
\hline
\rowcolor{lightgray!40} \textbf{Use Case} & \textbf{UC05: Verify Faculty Academic Credentials \& Tier Assignment} \\ \hline
\textbf{Goal} & System or administrator validates faculty academic credentials and assigns hierarchical trust badges (Tier 1--3). \\ \hline
\textbf{Preconditions} & Faculty account is registered and email verified; coordinator or admin has review privileges. \\ \hline
\textbf{Success End Condition} & Credential tier updated; verification badge displayed on expert profile. Message: \textit{``Expert verification tier updated to Tier-2.''} \\ \hline
\textbf{Failed End Condition} & Verification tier unchanged. Message: \textit{``Verification rejected: Inconclusive academic credentials.''} \\ \hline
\textbf{Primary Actors:} & Department Capstone Coordinator, System Administrator \\
\textbf{Secondary Actors:} & Faculty Advisor \\ \hline
\textbf{Trigger} & Faculty requests credential tier upgrade or coordinator reviews pending departmental faculty list. \\ \hline
\textbf{Description / Main Success Scenario} &
\begin{enumerate}[leftmargin=*]
    \item Coordinator opens the Faculty Verification queue on the administrative console.
    \item System displays faculty profile, department, institutional email domain, and synced publications.
    \item Coordinator inspects academic credentials and departmental appointment records.
    \item Coordinator selects verification tier: Tier-1 (Email Verified), Tier-2 (Scholar Verified), or Tier-3 (Department Head Validated).
    \item System persists updated tier in database and logs approval in audit trail.
    \item System sends in-app and email notification to faculty advisor confirming tier activation.
\end{enumerate} \\ \hline
\textbf{Alternative Flows} &
\begin{itemize}[leftmargin=*]
    \item[\textbf{3.a.}] \textbf{Credentials unverifiable or fraudulent:}
    \begin{itemize}[leftmargin=*]
        \item[3.a.1.] Coordinator rejects tier upgrade and inputs formal justification note.
        \item[3.a.2.] System marks account as unverified and notifies faculty member to provide proof.
    \end{itemize}
\end{itemize} \\ \hline
\textbf{Quality Requirements} & 
1. Restrict credential tier modification strictly to Coordinator and Administrator roles. \newline
2. Record reviewer identity, timestamp, and justification in immutable audit trail. \\ \hline
\end{tabularx}
\end{table}

\clearpage
\subsection{Use Case Description-06: Ingest Problem Statement \& Natural Language Parsing}
\label{subsec:uc06}

\begin{table}[H]
\centering
\small
\begin{tabularx}{\textwidth}{|l|X|}
\hline
\rowcolor{lightgray!40} \textbf{Use Case} & \textbf{UC06: Ingest Problem Statement \& Natural Language Parsing} \\ \hline
\textbf{Goal} & A student team inputs an unstructured project problem statement and initiates automated NLP processing. \\ \hline
\textbf{Preconditions} & Student lead has an active authenticated session; project proposal form is accessible. \\ \hline
\textbf{Success End Condition} & Problem statement stored; preliminary text tokens parsed and dispatched to entity extractor. Message: \textit{``Problem formulation ingested. Processing NLP extraction.''} \\ \hline
\textbf{Failed End Condition} & Ingestion rejected. Message: \textit{``Submission rejected: Problem statement must contain between 50 and 5,000 characters.''} \\ \hline
\textbf{Primary Actors:} & Student Project Lead \\
\textbf{Secondary Actors:} & None \\ \hline
\textbf{Trigger} & Student inputs problem text and clicks ``Analyze Proposal''. \\ \hline
\textbf{Description / Main Success Scenario} &
\begin{enumerate}[leftmargin=*]
    \item Student navigates to the ``Propose Capstone Problem'' portal.
    \item Student enters project title, target problem description, and envisioned methodology.
    \item System validates character count ($\ge 50$ and $\le 5,000$ characters).
    \item System cleanses input text, strips malicious scripts, and normalizes UTF-8 punctuation.
    \item System saves draft project entity in PostgreSQL.
    \item System automatically triggers Use Case UC07 (Extract Domain Taxonomy \& Skill Entities).
\end{enumerate} \\ \hline
\textbf{Alternative Flows} &
\begin{itemize}[leftmargin=*]
    \item[\textbf{3.a.}] \textbf{Text is under 50 characters (insufficient semantic context):}
    \begin{itemize}[leftmargin=*]
        \item[3.a.1.] System warns student: \textit{``Please provide a more detailed problem description to enable accurate AI matching.''}
    \end{itemize}
    \item[\textbf{3.b.}] \textbf{Text exceeds 5,000 characters:}
    \begin{itemize}[leftmargin=*]
        \item[3.b.1.] System displays character overflow warning and prevents submission until trimmed.
    \end{itemize}
\end{itemize} \\ \hline
\textbf{Quality Requirements} & 
1. Sanitize text against cross-site scripting (XSS) and SQL injection payloads. \newline
2. Ingest and save text within 200 milliseconds. \\ \hline
\end{tabularx}
\end{table}

\clearpage
\subsection{Use Case Description-07: Extract Domain Taxonomy \& Prerequisite Skill Entities}
\label{subsec:uc07}

\begin{table}[H]
\centering
\small
\begin{tabularx}{\textwidth}{|l|X|}
\hline
\rowcolor{lightgray!40} \textbf{Use Case} & \textbf{UC07: Extract Domain Taxonomy \& Prerequisite Skill Entities} \\ \hline
\textbf{Goal} & The AI inference pipeline automatically parses raw project text into macro academic domains and micro technical skill entities. \\ \hline
\textbf{Preconditions} & Problem statement successfully ingested via UC06; NLP microservice is operational. \\ \hline
\textbf{Success End Condition} & Domain tags and skill entity list generated with confidence scores. Message: \textit{``Taxonomy and skill entities extracted.''} \\ \hline
\textbf{Failed End Condition} & Extraction fails. Message: \textit{``NLP processing error: Unable to extract skills. Switching to keyword mode.''} \\ \hline
\textbf{Primary Actors:} & Student Project Lead \\
\textbf{Secondary Actors:} & None \\ \hline
\textbf{Trigger} & Automatically triggered following UC06 submission. \\ \hline
\textbf{Description / Main Success Scenario} &
\begin{enumerate}[leftmargin=*]
    \item NLP service receives normalized problem statement text.
    \item Service executes zero-shot transformer classification against academic discipline taxonomy.
    \item Service identifies top 3 academic domains (e.g., \texttt{Embedded Systems}, \texttt{Machine Learning}) with confidence $\ge 0.65$.
    \item Service executes Named Entity Recognition (NER) pipeline to extract technical competencies, libraries, and hardware platforms.
    \item Service maps extracted entities to canonical taxonomy terms (e.g., mapping ``Tensorflow 2'' $\rightarrow$ \texttt{TensorFlow}).
    \item System structures extracted attributes into JSON payload and transitions to UC08.
\end{enumerate} \\ \hline
\textbf{Alternative Flows} &
\begin{itemize}[leftmargin=*]
    \item[\textbf{3.a.}] \textbf{No domain meets 0.65 confidence threshold:}
    \begin{itemize}[leftmargin=*]
        \item[3.a.1.] System selects top 2 nearest domains and flags them as tentative for student confirmation.
    \end{itemize}
    \item[\textbf{4.a.}] \textbf{NLP model timeout or container crash:}
    \begin{itemize}[leftmargin=*]
        \item[4.a.1.] System falls back to regex-based keyword matching against academic skill catalog.
        \item[4.a.2.] System logs fallback event in administrative monitoring console.
    \end{itemize}
\end{itemize} \\ \hline
\textbf{Quality Requirements} & 
1. Achieve $\ge 85\%$ precision on domain classification across academic benchmark datasets. \newline
2. Complete zero-shot classification and entity extraction within 1.5 seconds. \\ \hline
\end{tabularx}
\end{table}

\clearpage
\subsection{Use Case Description-08: Refine Extracted Requirements \& Skill Weights}
\label{subsec:uc08}

\begin{table}[H]
\centering
\small
\begin{tabularx}{\textwidth}{|l|X|}
\hline
\rowcolor{lightgray!40} \textbf{Use Case} & \textbf{UC08: Refine Extracted Requirements \& Skill Weights} \\ \hline
\textbf{Goal} & The student team reviews, adjusts, adds, removes, or re-weights AI-extracted domains and technical skill prerequisites. \\ \hline
\textbf{Preconditions} & Entities extracted via UC07; review card displayed on student screen. \\ \hline
\textbf{Success End Condition} & Validated project requirement vector saved and ready for vector search. Message: \textit{``Requirements confirmed. Searching expert database.''} \\ \hline
\textbf{Failed End Condition} & Refinement cancelled. Message: \textit{``Changes discarded.''} \\ \hline
\textbf{Primary Actors:} & Student Project Lead \\
\textbf{Secondary Actors:} & None \\ \hline
\textbf{Trigger} & Student reviews the interactive entity tags generated by the NLP pipeline. \\ \hline
\textbf{Description / Main Success Scenario} &
\begin{enumerate}[leftmargin=*]
    \item Student inspects extracted academic domains and technical skill tags.
    \item Student adjusts importance sliders for specific skills (e.g., setting \texttt{PyTorch} to High [1.0] and \texttt{Docker} to Low [0.3]).
    \item Student types to add missing specialized skill tags using autocomplete suggestions.
    \item Student clicks ``Dismiss'' on any irrelevant entities extracted by the NLP model.
    \item Student clicks ``Confirm Requirements \& Find Advisors''.
    \item System persists finalized requirement weights and triggers vector matchmaking (UC09).
\end{enumerate} \\ \hline
\textbf{Alternative Flows} &
\begin{itemize}[leftmargin=*]
    \item[\textbf{3.a.}] \textbf{Student enters custom skill not in system taxonomy:}
    \begin{itemize}[leftmargin=*]
        \item[3.a.1.] System accepts custom entity, flags it for administrator taxonomy review, and indexes it.
    \end{itemize}
    \item[\textbf{5.a.}] \textbf{All extracted skills removed by student:}
    \begin{itemize}[leftmargin=*]
        \item[5.a.1.] System prompts student: \textit{``At least one technical skill is required for matching. Please retain or add skills.''}
    \end{itemize}
\end{itemize} \\ \hline
\textbf{Quality Requirements} & 
1. Autocomplete response time under 100 milliseconds across 2,000+ taxonomy terms. \newline
2. Immediate reactive UI updates when toggling skill sliders. \\ \hline
\end{tabularx}
\end{table}

\clearpage
\subsection{Use Case Description-09: Execute Dense Semantic Vector Search}
\label{subsec:uc09}

\begin{table}[H]
\centering
\small
\begin{tabularx}{\textwidth}{|l|X|}
\hline
\rowcolor{lightgray!40} \textbf{Use Case} & \textbf{UC09: Execute Dense Semantic Vector Search} \\ \hline
\textbf{Goal} & Retrieve nearest candidate faculty and mentor profiles from high-dimensional vector space using cosine distance. \\ \hline
\textbf{Preconditions} & Project requirements confirmed; Qdrant vector database is reachable and populated. \\ \hline
\textbf{Success End Condition} & Top-50 nearest candidate vectors retrieved with cosine similarity scores. Message: \textit{``Candidate vector retrieval completed.''} \\ \hline
\textbf{Failed End Condition} & Search fails. Message: \textit{``Vector search unavailable: Reverting to exact-match database retrieval.''} \\ \hline
\textbf{Primary Actors:} & Student Project Lead \\
\textbf{Secondary Actors:} & System Administrator \\ \hline
\textbf{Trigger} & Student triggers advisor discovery following requirement confirmation. \\ \hline
\textbf{Description / Main Success Scenario} &
\begin{enumerate}[leftmargin=*]
    \item Search service takes confirmed problem text and weighted skill tokens.
    \item Service encodes input text into a 384-dimensional dense vector using transformer embedding worker.
    \item Service connects to Qdrant vector store and initiates approximate nearest neighbor (ANN) search.
    \item HNSW index evaluates cosine distance against indexed faculty research corpora:
    $$\cos(\vec{q}, \vec{d}_i) = \frac{\vec{q} \cdot \vec{d}_i}{\|\vec{q}\| \|\vec{d}_i\|}$$
    \item Qdrant returns top-50 candidate expert IDs with vector similarity scores ($S_{\text{dense}} \in [0, 1]$).
    \item Service passes candidate IDs and scores to Hybrid Match Ranking Engine (UC10).
\end{enumerate} \\ \hline
\textbf{Alternative Flows} &
\begin{itemize}[leftmargin=*]
    \item[\textbf{3.a.}] \textbf{Qdrant cluster connection timeout:}
    \begin{itemize}[leftmargin=*]
        \item[3.a.1.] System catches network timeout and initiates fallback to PostgreSQL trigram keyword search.
        \item[3.a.2.] System delivers fallback results with notification: \textit{``Displaying keyword-matched results.''}
    \end{itemize}
\end{itemize} \\ \hline
\textbf{Quality Requirements} & 
1. Sub-200ms p95 latency for ANN vector search across 10,000 indexed faculty profiles. \newline
2. Seamless graceful degradation to keyword fallback if vector cluster is offline. \\ \hline
\end{tabularx}
\end{table}

\clearpage
\subsection{Use Case Description-10: Multi-Criteria Hybrid Match Ranking \& Scoring}
\label{subsec:uc10}

\begin{table}[H]
\centering
\small
\begin{tabularx}{\textwidth}{|l|X|}
\hline
\rowcolor{lightgray!40} \textbf{Use Case} & \textbf{UC10: Multi-Criteria Hybrid Match Ranking \& Scoring} \\ \hline
\textbf{Goal} & Compute composite matching scores combining semantic similarity, skill overlap, publication recency, and availability. \\ \hline
\textbf{Preconditions} & Candidate vector scores returned from UC09; relational profile attributes loaded from PostgreSQL. \\ \hline
\textbf{Success End Condition} & Ranked list of candidate experts generated with composite score $S_{\text{hybrid}} \in [0, 100]$. Message: \textit{``Match ranking completed.''} \\ \hline
\textbf{Failed End Condition} & Scoring failed. Message: \textit{``Unable to compute rankings; please retry.''} \\ \hline
\textbf{Primary Actors:} & Student Project Lead \\
\textbf{Secondary Actors:} & Faculty Advisor \\ \hline
\textbf{Trigger} & Candidate vector retrieval completes. \\ \hline
\textbf{Description / Main Success Scenario} &
\begin{enumerate}[leftmargin=*]
    \item Ranking engine fetches full metadata for top-50 candidate IDs from PostgreSQL.
    \item Engine calculates normalized skill overlap ratio:
    $$S_{\text{skill}} = \frac{|\mathcal{K}_{\text{req}} \cap \mathcal{K}_{\text{expert}}|}{|\mathcal{K}_{\text{req}}|}$$
    \item Engine calculates publication recency and domain citation weight ($S_{\text{pub}}$).
    \item Engine applies verification tier multiplier ($S_{\text{tier}} \in \{1.0, 1.15, 1.25\}$).
    \item Engine evaluates active advisory quota capacity multiplier ($S_{\text{avail}}$).
    \item Engine computes overall hybrid composite score:
    $$S_{\text{hybrid}} = \left(\omega_1 S_{\text{dense}} + \omega_2 S_{\text{skill}} + \omega_3 S_{\text{pub}} + \omega_4 S_{\text{tier}}\right) \times S_{\text{avail}}$$
    \item Engine sorts candidates in descending score order and delivers results to client.
\end{enumerate} \\ \hline
\textbf{Alternative Flows} &
\begin{itemize}[leftmargin=*]
    \item[\textbf{5.a.}] \textbf{Advisor has 0 remaining supervisory slots ($S_{\text{avail}} = 0$ for primary supervision):}
    \begin{itemize}[leftmargin=*]
        \item[5.a.1.] System adjusts profile status to ``Quota Full'' and flags candidate for co-advisory only.
    \end{itemize}
\end{itemize} \\ \hline
\textbf{Quality Requirements} & 
1. Calculate composite scores and rank 50 candidates in under 300 milliseconds. \newline
2. Composite score formulation normalized strictly between 0 and 100. \\ \hline
\end{tabularx}
\end{table}

\clearpage
\subsection{Use Case Description-11: Generate Explainable AI (XAI) Recommendation Breakdown}
\label{subsec:uc11}

\begin{table}[H]
\centering
\small
\begin{tabularx}{\textwidth}{|l|X|}
\hline
\rowcolor{lightgray!40} \textbf{Use Case} & \textbf{UC11: Generate Explainable AI (XAI) Recommendation Breakdown} \\ \hline
\textbf{Goal} & Provide transparent, human-interpretable justifications for each recommended candidate expert. \\ \hline
\textbf{Preconditions} & Candidate expert list ranked via UC10; student selects a candidate card. \\ \hline
\textbf{Success End Condition} & Detailed XAI justification card and percentage attribution breakdown rendered. Message: \textit{``Match justification loaded.''} \\ \hline
\textbf{Failed End Condition} & Justification unavailable. Message: \textit{``Detailed explanation could not be generated.''} \\ \hline
\textbf{Primary Actors:} & Student Project Lead \\
\textbf{Secondary Actors:} & Faculty Advisor, Department Coordinator \\ \hline
\textbf{Trigger} & Student clicks ``Why this match?'' on a candidate expert profile card. \\ \hline
\textbf{Description / Main Success Scenario} &
\begin{enumerate}[leftmargin=*]
    \item Student clicks ``Why this match?'' on Dr. Farhana Rahman's profile card.
    \item System retrieves score component vectors computed during UC10.
    \item System formats component percentage contributions into an interactive breakdown modal:
    \begin{itemize}[leftmargin=*]
        \item 45\% Semantic Vector Similarity (highlighting 3 matched publication abstracts).
        \item 30\% Prerequisite Skill Overlap (displaying matched skills in green).
        \item 15\% Verified Academic Credential Tier (Tier-2 Scholar Verified).
        \item 10\% Supervisory Bandwidth (2 active slots available).
    \end{itemize}
    \item System renders a radar chart comparing student project requirements against expert proficiencies.
    \item Student inspects justification and gains confidence in candidate suitability.
\end{enumerate} \\ \hline
\textbf{Alternative Flows} &
\begin{itemize}[leftmargin=*]
    \item[\textbf{3.a.}] \textbf{Match is driven almost exclusively by keyword overlap rather than semantic papers:}
    \begin{itemize}[leftmargin=*]
        \item[3.a.1.] System displays an advisory hint: \textit{``Strong practical toolchain overlap; limited recent publication history in this exact domain.''}
    \end{itemize}
\end{itemize} \\ \hline
\textbf{Quality Requirements} & 
1. Render XAI visual drawer and radar chart within 200 milliseconds of user click. \newline
2. All percentage contributions must sum precisely to 100\%. \\ \hline
\end{tabularx}
\end{table}

\clearpage
\subsection{Use Case Description-12: Filter Candidates by Department, Quota \& Availability}
\label{subsec:uc12}

\begin{table}[H]
\centering
\small
\begin{tabularx}{\textwidth}{|l|X|}
\hline
\rowcolor{lightgray!40} \textbf{Use Case} & \textbf{UC12: Filter Candidates by Department, Quota \& Availability} \\ \hline
\textbf{Goal} & Allow students to dynamically refine candidate results using faceted academic criteria. \\ \hline
\textbf{Preconditions} & Candidate recommendation results active on the search screen. \\ \hline
\textbf{Success End Condition} & Filtered candidate list updated in real time. Message: \textit{``Filters applied.''} \\ \hline
\textbf{Failed End Condition} & No candidates match filter criteria. Message: \textit{``No experts match the selected filters. Try broadening your criteria.''} \\ \hline
\textbf{Primary Actors:} & Student Project Lead \\
\textbf{Secondary Actors:} & None \\ \hline
\textbf{Trigger} & Student selects filter checkboxes on the search sidebar. \\ \hline
\textbf{Description / Main Success Scenario} &
\begin{enumerate}[leftmargin=*]
    \item Student opens the faceted search filter panel.
    \item Student selects filter criteria: Department = ``Software Engineering'', Quota Status = ``Available Now'', Credential Tier = ``Tier-2+''.
    \item Client dispatches filter parameters to backend API.
    \item System applies filter predicates against cached recommendation candidate set.
    \item System updates UI with filtered candidate list and displays matching count.
\end{enumerate} \\ \hline
\textbf{Alternative Flows} &
\begin{itemize}[leftmargin=*]
    \item[\textbf{4.a.}] \textbf{Filter results in 0 matches:}
    \begin{itemize}[leftmargin=*]
        \item[4.a.1.] System displays ``Zero Matches Found'' empty state card.
        \item[4.a.2.] System provides a 1-click ``Reset All Filters'' action to restore recommendations.
    \end{itemize}
\end{itemize} \\ \hline
\textbf{Quality Requirements} & 
1. Apply faceted filters and update client view in under 150 milliseconds. \newline
2. Maintain active filter state across browser tab refreshes. \\ \hline
\end{tabularx}
\end{table}

\clearpage
\subsection{Use Case Description-13: Perform Side-by-Side Expert Comparison}
\label{subsec:uc13}

\begin{table}[H]
\centering
\small
\begin{tabularx}{\textwidth}{|l|X|}
\hline
\rowcolor{lightgray!40} \textbf{Use Case} & \textbf{UC13: Perform Side-by-Side Expert Comparison} \\ \hline
\textbf{Goal} & Compare up to three candidate advisors across standardized academic and availability metrics. \\ \hline
\textbf{Preconditions} & Student has identified multiple candidate experts in search results. \\ \hline
\textbf{Success End Condition} & Side-by-side comparison matrix rendered with differential metric highlights. Message: \textit{``Comparison matrix ready.''} \\ \hline
\textbf{Failed End Condition} & Comparison failed. Message: \textit{``Unable to load comparison: Select between 2 and 3 candidates.''} \\ \hline
\textbf{Primary Actors:} & Student Project Lead \\
\textbf{Secondary Actors:} & None \\ \hline
\textbf{Trigger} & Student checks ``Compare'' on two or three candidate cards and clicks ``Compare Selected''. \\ \hline
\textbf{Description / Main Success Scenario} &
\begin{enumerate}[leftmargin=*]
    \item Student selects two or three candidate experts from search results.
    \item Student clicks ``Compare Selected Experts''.
    \item System loads full profile attributes for chosen experts from database.
    \item System renders a tabular comparative matrix aligning:
    \begin{itemize}[leftmargin=*]
        \item Overall Match Score and Domain Similarity.
        \item Matched vs. Missing Technical Skills.
        \item Total Citations, h-index, and Recent Publication Count (last 3 years).
        \item Available Supervisory Capacity and Office Hour Availability.
        \item Average Post-Collaboration Student Rating.
    \end{itemize}
    \item Student analyzes relative strengths and selects preferred primary supervisor.
\end{enumerate} \\ \hline
\textbf{Alternative Flows} &
\begin{itemize}[leftmargin=*]
    \item[\textbf{1.a.}] \textbf{Student attempts to select more than 3 candidates:}
    \begin{itemize}[leftmargin=*]
        \item[1.a.1.] System displays notification: \textit{``You can compare a maximum of 3 candidates simultaneously.''}
    \end{itemize}
\end{itemize} \\ \hline
\textbf{Quality Requirements} & 
1. Render comparison matrix in under 300 milliseconds. \newline
2. Highlight differential metrics clearly with color-coded comparison badges. \\ \hline
\end{tabularx}
\end{table}

\clearpage
\subsection{Use Case Description-14: Analyze Multidisciplinary Team Skill Gap ($\Delta \mathcal{K}$)}
\label{subsec:uc14}

\begin{table}[H]
\centering
\small
\begin{tabularx}{\textwidth}{|l|X|}
\hline
\rowcolor{lightgray!40} \textbf{Use Case} & \textbf{UC14: Analyze Multidisciplinary Team Skill Gap ($\Delta \mathcal{K}$)} \\ \hline
\textbf{Goal} & Detect unaddressed prerequisite technical skills in a student project team and evaluate project viability. \\ \hline
\textbf{Preconditions} & Project proposal requirements defined; student team member profiles linked to project. \\ \hline
\textbf{Success End Condition} & Skill gap report generated showing addressed vs. missing prerequisite competencies. Message: \textit{``Skill gap analysis complete.''} \\ \hline
\textbf{Failed End Condition} & Analysis failed. Message: \textit{``Unable to compute skill gaps: Team roster missing.''} \\ \hline
\textbf{Primary Actors:} & Student Project Lead, Department Capstone Coordinator \\
\textbf{Secondary Actors:} & None \\ \hline
\textbf{Trigger} & Student lead or coordinator clicks ``Analyze Team Skill Gaps''. \\ \hline
\textbf{Description / Main Success Scenario} &
\begin{enumerate}[leftmargin=*]
    \item System retrieves project prerequisite skill set $\mathcal{K}_{\text{req}}$ from project record.
    \item System queries profiles of all confirmed team members to compute collective skill union:
    $$\mathcal{K}_{\text{team}} = \bigcup_{m \in \text{Members}} \mathcal{K}_m$$
    \item System calculates the unaddressed deficit set: $\Delta \mathcal{K} = \mathcal{K}_{\text{req}} \setminus \mathcal{K}_{\text{team}}$.
    \item System computes team readiness score:
    $$R_{\text{team}} = 1 - \frac{|\Delta \mathcal{K}|}{|\mathcal{K}_{\text{req}}|}$$
    \item System displays visual breakdown highlighting covered skills (green) and critical deficits (red).
    \item System displays prompt: \textit{``Would you like to discover co-advisors or peers to resolve missing skills?''}
\end{enumerate} \\ \hline
\textbf{Alternative Flows} &
\begin{itemize}[leftmargin=*]
    \item[\textbf{3.a.}] \textbf{Deficit set is empty ($\Delta \mathcal{K} = \emptyset$):}
    \begin{itemize}[leftmargin=*]
        \item[3.a.1.] System displays ``100\% Team Skill Coverage'' green banner and bypasses co-advisor prompt.
    \end{itemize}
\end{itemize} \\ \hline
\textbf{Quality Requirements} & 
1. Compute skill set union and set difference in under 50 milliseconds. \newline
2. Accurately flag interdisciplinary skill deficits across department boundaries. \\ \hline
\end{tabularx}
\end{table}

\clearpage
\subsection{Use Case Description-15: Run Set-Cover Optimization for Co-Advisor Selection}
\label{subsec:uc15}

\begin{table}[H]
\centering
\small
\begin{tabularx}{\textwidth}{|l|X|}
\hline
\rowcolor{lightgray!40} \textbf{Use Case} & \textbf{UC15: Run Set-Cover Optimization for Co-Advisor Selection} \\ \hline
\textbf{Goal} & Formulate and solve a combinatorial set-cover problem to recommend the minimal set of co-advisors to cover team gaps. \\ \hline
\textbf{Preconditions} & Skill deficit set $\Delta \mathcal{K} \ne \emptyset$ identified via UC14; faculty pool indexed. \\ \hline
\textbf{Success End Condition} & Optimal co-advisor recommendations presented with simulated post-addition coverage. Message: \textit{``Optimal co-advisor candidates identified.''} \\ \hline
\textbf{Failed End Condition} & Optimization failed. Message: \textit{``No co-advisors found covering the specified skill gaps.''} \\ \hline
\textbf{Primary Actors:} & Student Project Lead, Department Capstone Coordinator \\
\textbf{Secondary Actors:} & Faculty Advisor \\ \hline
\textbf{Trigger} & User selects ``Find Co-Advisors to Bridge Skill Gaps''. \\ \hline
\textbf{Description / Main Success Scenario} &
\begin{enumerate}[leftmargin=*]
    \item System loads unaddressed skill set $\Delta \mathcal{K}$.
    \item System queries candidate faculty members outside the primary supervisor's department.
    \item System executes a greedy approximation algorithm for weighted set-cover to select candidates that maximize:
    $$\arg\max_{e \in \text{Experts}} \frac{|\mathcal{K}_e \cap \Delta \mathcal{K}|}{\text{Cost}(e)}$$
    \item System ranks candidate co-advisors based on marginal skill contribution and quota availability.
    \item System presents top co-advisor recommendations with a simulation card: \textit{``Adding Dr. Tanvir Ahmed (Bioinformatics) increases team coverage from 40\% to 92\%.''}
    \item Student lead clicks ``Propose Co-Advisory Collaboration''.
\end{enumerate} \\ \hline
\textbf{Alternative Flows} &
\begin{itemize}[leftmargin=*]
    \item[\textbf{3.a.}] \textbf{No single faculty covers all deficits:}
    \begin{itemize}[leftmargin=*]
        \item[3.a.1.] System recommends a paired combination of one co-advisor and one external industry mentor.
    \end{itemize}
\end{itemize} \\ \hline
\textbf{Quality Requirements} & 
1. Execute greedy set-cover algorithm in under 200 milliseconds. \newline
2. Guarantee that recommended set strictly reduces cardinality of $\Delta \mathcal{K}$. \\ \hline
\end{tabularx}
\end{table}

\clearpage
\subsection{Use Case Description-16: Submit Collaboration Advisory Request}
\label{subsec:uc16}

\begin{table}[H]
\centering
\small
\begin{tabularx}{\textwidth}{|l|X|}
\hline
\rowcolor{lightgray!40} \textbf{Use Case} & \textbf{UC16: Submit Collaboration Advisory Request} \\ \hline
\textbf{Goal} & A student project team submits a formal, structured capstone advisory proposal to a selected faculty expert. \\ \hline
\textbf{Preconditions} & Student lead has active project draft; selected faculty has available quota ($>0$ active slots). \\ \hline
\textbf{Success End Condition} & Advisory proposal record created; FSM transitions to \texttt{SUBMITTED}; notification sent. Message: \textit{``Advisory request submitted successfully.''} \\ \hline
\textbf{Failed End Condition} & Submission blocked. Message: \textit{``Unable to submit request: Faculty quota full or pending request already exists.''} \\ \hline
\textbf{Primary Actors:} & Student Project Lead \\
\textbf{Secondary Actors:} & Faculty Advisor \\ \hline
\textbf{Trigger} & Student clicks ``Send Advisory Proposal'' on expert profile page. \\ \hline
\textbf{Description / Main Success Scenario} &
\begin{enumerate}[leftmargin=*]
    \item Student opens the advisory proposal modal.
    \item Student reviews pre-populated project synopsis and AI alignment score.
    \item Student enters structured proposal items: envisioned milestones, expected meeting frequency, and personal motivation statement.
    \item Student selects advisory scope: Primary Supervisor or Co-Advisor.
    \item System validates that the student team has no other pending requests to the same professor.
    \item System verifies that the target faculty member's available quota is $\ge 1$.
    \item System creates an \texttt{AdvisoryRequest} record in PostgreSQL with status \texttt{SUBMITTED}.
    \item System starts a 7-day auto-response countdown timer.
    \item System delivers dashboard notification and email digest to the faculty advisor.
\end{enumerate} \\ \hline
\textbf{Alternative Flows} &
\begin{itemize}[leftmargin=*]
    \item[\textbf{6.a.}] \textbf{Faculty quota reached while student was editing form:}
    \begin{itemize}[leftmargin=*]
        \item[6.a.1.] System blocks submission: \textit{``Professor's quota was filled by another group. Proposal cannot be submitted.''}
        \item[6.a.2.] System prompts student to select an alternative recommended advisor.
    \end{itemize}
\end{itemize} \\ \hline
\textbf{Quality Requirements} & 
1. Enforce atomic database transaction to prevent race conditions on quota capacity. \newline
2. Dispatch confirmation notification within 5 seconds of submission. \\ \hline
\end{tabularx}
\end{table}

\clearpage
\subsection{Use Case Description-17: Manage Advisory Request Lifecycle (Accept / Decline / Negotiate)}
\label{subsec:uc17}

\begin{table}[H]
\centering
\small
\begin{tabularx}{\textwidth}{|l|X|}
\hline
\rowcolor{lightgray!40} \textbf{Use Case} & \textbf{UC17: Manage Advisory Request Lifecycle (Accept / Decline / Negotiate)} \\ \hline
\textbf{Goal} & A faculty advisor evaluates an incoming student collaboration proposal and executes an FSM state transition. \\ \hline
\textbf{Preconditions} & Faculty advisor is signed in; advisory proposal in \texttt{SUBMITTED} or \texttt{UNDER\_REVIEW} state. \\ \hline
\textbf{Success End Condition} & FSM state updated to \texttt{ACCEPTED}, \texttt{DECLINED}, or \texttt{NEGOTIATION}. Message: \textit{``Advisory request status updated.''} \\ \hline
\textbf{Failed End Condition} & State transition rejected. Message: \textit{``Invalid state transition or proposal expired.''} \\ \hline
\textbf{Primary Actors:} & Faculty Advisor \\
\textbf{Secondary Actors:} & Student Project Lead, Department Capstone Coordinator \\ \hline
\textbf{Trigger} & Faculty advisor opens a pending request from their ``Incoming Solicitations'' queue. \\ \hline
\textbf{Description / Main Success Scenario} &
\begin{enumerate}[leftmargin=*]
    \item Faculty advisor opens pending proposal in the advisory portal.
    \item System updates proposal state to \texttt{UNDER\_REVIEW}.
    \item Advisor reviews student project statement, proposed milestones, and XAI match score.
    \item Advisor selects action:
    \begin{itemize}[leftmargin=*]
        \item \textbf{Accept:} FSM transitions to \texttt{ACCEPTED}; advisor quota decremented; active workspace created.
        \item \textbf{Negotiate:} Advisor enters milestone modification feedback; FSM transitions to \texttt{NEGOTIATION}.
        \item \textbf{Decline:} Advisor inputs polite declination reason; FSM transitions to \texttt{DECLINED}.
    \end{itemize}
    \item System persists state transition with timestamp and user ID in audit trail.
    \item System triggers real-time notification to student project team.
\end{enumerate} \\ \hline
\textbf{Alternative Flows} &
\begin{itemize}[leftmargin=*]
    \item[\textbf{4.a.}] \textbf{7-day response timer expires without advisor action:}
    \begin{itemize}[leftmargin=*]
        \item[4.a.1.] System background cron job automatically transitions status to \texttt{EXPIRED}.
        \item[4.a.2.] System releases student team from waiting lock and suggests next-best ranked advisors.
    \end{itemize}
\end{itemize} \\ \hline
\textbf{Quality Requirements} & 
1. Enforce strict mathematical state transition rules with no illegal skips. \newline
2. Log every transition and feedback comment in immutable audit ledger. \\ \hline
\end{tabularx}
\end{table}

\clearpage
\subsection{Use Case Description-18: Enforce Faculty Advisory Quotas \& Blackout Windows}
\label{subsec:uc18}

\begin{table}[H]
\centering
\small
\begin{tabularx}{\textwidth}{|l|X|}
\hline
\rowcolor{lightgray!40} \textbf{Use Case} & \textbf{UC18: Enforce Faculty Advisory Quotas \& Blackout Windows} \\ \hline
\textbf{Goal} & Maintain balanced faculty supervision loads and prevent inbox burnout by enforcing hard quota caps and blackout dates. \\ \hline
\textbf{Preconditions} & Faculty member has active profile; department capstone quota policies defined. \\ \hline
\textbf{Success End Condition} & Quota settings updated; availability status dynamically toggled. Message: \textit{``Supervisory quota settings saved.''} \\ \hline
\textbf{Failed End Condition} & Settings rejected. Message: \textit{``Invalid quota value: Must be between 1 and 10 groups.''} \\ \hline
\textbf{Primary Actors:} & Faculty Advisor, Department Capstone Coordinator \\
\textbf{Secondary Actors:} & System Administrator \\ \hline
\textbf{Trigger} & Faculty advisor configures settings or an advisory proposal is accepted. \\ \hline
\textbf{Description / Main Success Scenario} &
\begin{enumerate}[leftmargin=*]
    \item Faculty advisor accesses the ``Advisory Capacity \& Quotas'' settings tab.
    \item Advisor sets maximum active capstone supervision groups (e.g., 4 groups).
    \item Advisor designates blackout calendar windows (e.g., sabbatical dates or conference travel).
    \item System validates settings against departmental statutory policy limits.
    \item System recalculates current active count: $\text{ActiveGroups} = \text{Count}(\text{AcceptedProposals})$.
    \item If $\text{ActiveGroups} \ge \text{MaxQuota}$, system automatically transitions status to \texttt{QUOTA\_REACHED}.
    \item System updates expert availability index in vector store to depress primary supervisor ranking.
\end{enumerate} \\ \hline
\textbf{Alternative Flows} &
\begin{itemize}[leftmargin=*]
    \item[\textbf{4.a.}] \textbf{Coordinator enacts departmental quota override:}
    \begin{itemize}[leftmargin=*]
        \item[4.a.1.] Coordinator elevates faculty cap by +1 group due to special research necessity.
        \item[4.a.2.] System records coordinator override rationale in audit log.
    \end{itemize}
\end{itemize} \\ \hline
\textbf{Quality Requirements} & 
1. Real-time quota enforcement with atomic database locking on acceptance. \newline
2. Asynchronous status update to vector search index within 2 seconds. \\ \hline
\end{tabularx}
\end{table}

\clearpage
\subsection{Use Case Description-19: Schedule Mentorship Consultation Session}
\label{subsec:uc19}

\begin{table}[H]
\centering
\small
\begin{tabularx}{\textwidth}{|l|X|}
\hline
\rowcolor{lightgray!40} \textbf{Use Case} & \textbf{UC19: Schedule Mentorship Consultation Session} \\ \hline
\textbf{Goal} & An approved student team schedules an office-hour consultation session with their assigned advisor. \\ \hline
\textbf{Preconditions} & Active accepted advisory relationship exists; advisor has published open consultation slots. \\ \hline
\textbf{Success End Condition} & Consultation booked; calendar invite generated; status \texttt{SCHEDULED}. Message: \textit{``Consultation booked successfully.''} \\ \hline
\textbf{Failed End Condition} & Booking rejected. Message: \textit{``Slot no longer available or agenda missing.''} \\ \hline
\textbf{Primary Actors:} & Student Project Lead \\
\textbf{Secondary Actors:} & Faculty Advisor, Industry Mentor \\ \hline
\textbf{Trigger} & Student opens the ``Schedule Consultation'' view in active project workspace. \\ \hline
\textbf{Description / Main Success Scenario} &
\begin{enumerate}[leftmargin=*]
    \item Student opens the advisor's consultation calendar.
    \item System displays open 30-minute office-hour slots for the upcoming two weeks.
    \item Student selects an open slot (e.g., Thursday 3:00 PM -- 3:30 PM).
    \item Student enters mandatory session agenda and attaches discussion documents.
    \item System verifies that the slot has not been concurrently reserved.
    \item System creates consultation record in database and marks slot as \texttt{BOOKED}.
    \item System generates an iCalendar (.ics) invite and sends email confirmation to both parties.
    \item System logs scheduled session in the project workspace milestone tracker.
\end{enumerate} \\ \hline
\textbf{Alternative Flows} &
\begin{itemize}[leftmargin=*]
    \item[\textbf{5.a.}] \textbf{Concurrent booking conflict (another student books slot simultaneously):}
    \begin{itemize}[leftmargin=*]
        \item[5.a.1.] System detects database slot lock collision.
        \item[5.a.2.] System alerts student: \textit{``Slot was just booked by another team. Please choose another time.''}
    \end{itemize}
\end{itemize} \\ \hline
\textbf{Quality Requirements} & 
1. Prevent double-booking using strict PostgreSQL row-level locks on calendar slots. \newline
2. Dispatch calendar invitations and notifications within 3 seconds of booking. \\ \hline
\end{tabularx}
\end{table}

\clearpage
\subsection{Use Case Description-20: Exchange Contextual Messages \& Project Milestones}
\label{subsec:uc20}

\begin{table}[H]
\centering
\small
\begin{tabularx}{\textwidth}{|l|X|}
\hline
\rowcolor{lightgray!40} \textbf{Use Case} & \textbf{UC20: Exchange Contextual Messages \& Project Milestones} \\ \hline
\textbf{Goal} & Students and advisors communicate, share milestone reports, and log technical feedback inside project workspace. \\ \hline
\textbf{Preconditions} & User belongs to an active \texttt{ACCEPTED} project workspace. \\ \hline
\textbf{Success End Condition} & Message/file persisted; conversation stream updated in real time. Message: \textit{``Message posted.''} \\ \hline
\textbf{Failed End Condition} & Transmission failed. Message: \textit{``Failed to send message: File type not permitted or network offline.''} \\ \hline
\textbf{Primary Actors:} & Student Project Lead, Faculty Advisor, Industry Mentor \\
\textbf{Secondary Actors:} & None \\ \hline
\textbf{Trigger} & User inputs text or attaches a document in the workspace discussion channel. \\ \hline
\textbf{Description / Main Success Scenario} &
\begin{enumerate}[leftmargin=*]
    \item User opens the workspace discussion channel for their capstone project.
    \item User enters message text or selects a milestone artifact (e.g., \texttt{Architecture\_v1.pdf}).
    \item System validates file attachment format (PDF, DOCX, ZIP; max 25 MB).
    \item System scans file for malicious payloads and stores it in encrypted object storage.
    \item System saves message record to database linked to the project and author ID.
    \item System broadcasts new message via WebSocket connection to active workspace participants.
    \item System updates unread message badge count for offline team members.
\end{enumerate} \\ \hline
\textbf{Alternative Flows} &
\begin{itemize}[leftmargin=*]
    \item[\textbf{3.a.}] \textbf{Attachment exceeds 25 MB limit:}
    \begin{itemize}[leftmargin=*]
        \item[3.a.1.] System rejects file and suggests providing an external repository link (e.g., GitHub/Drive).
    \end{itemize}
\end{itemize} \\ \hline
\textbf{Quality Requirements} & 
1. WebSocket message delivery latency under 100 milliseconds for connected clients. \newline
2. All messages encrypted at rest using AES-256-GCM. \\ \hline
\end{tabularx}
\end{table}

\clearpage
\subsection{Use Case Description-21: Submit Post-Collaboration Evaluation \& Endorsement}
\label{subsec:uc21}

\begin{table}[H]
\centering
\small
\begin{tabularx}{\textwidth}{|l|X|}
\hline
\rowcolor{lightgray!40} \textbf{Use Case} & \textbf{UC21: Submit Post-Collaboration Evaluation \& Endorsement} \\ \hline
\textbf{Goal} & Record mutual structured evaluations, quantitative ratings, and verified skill endorsements upon capstone completion. \\ \hline
\textbf{Preconditions} & Capstone project marked as completed by coordinator; advisory relationship closed. \\ \hline
\textbf{Success End Condition} & Mutual evaluation saved; reputation scores updated; skill attestations verified. Message: \textit{``Evaluation submitted successfully.''} \\ \hline
\textbf{Failed End Condition} & Submission rejected. Message: \textit{``Evaluation submission incomplete: Required rubric items missing.''} \\ \hline
\textbf{Primary Actors:} & Student Project Lead, Faculty Advisor, Industry Mentor \\
\textbf{Secondary Actors:} & Department Capstone Coordinator \\ \hline
\textbf{Trigger} & System prompts user for post-project review upon semester conclusion. \\ \hline
\textbf{Description / Main Success Scenario} &
\begin{enumerate}[leftmargin=*]
    \item User accesses the ``Project Completion Review'' portal.
    \item User completes a structured quantitative rubric (responsiveness, technical depth, milestone adherence).
    \item User inputs qualitative feedback comments.
    \item Advisor endorses specific skills demonstrated by students during the project.
    \item System validates that all mandatory evaluation fields are completed.
    \item System records evaluation in PostgreSQL with status \texttt{COMMITTED}.
    \item System updates the recipient's aggregate reputation score and attaches skill badges.
\end{enumerate} \\ \hline
\textbf{Alternative Flows} &
\begin{itemize}[leftmargin=*]
    \item[\textbf{3.a.}] \textbf{User submits derogatory or non-academic commentary:}
    \begin{itemize}[leftmargin=*]
        \item[3.a.1.] Automated content moderation filter flags text and routes review to Department Coordinator.
    \end{itemize}
\end{itemize} \\ \hline
\textbf{Quality Requirements} & 
1. Maintain confidentiality of student reviews until both parties submit or deadline passes. \newline
2. Update aggregated reputation metric immediately upon submission commit. \\ \hline
\end{tabularx}
\end{table}

\clearpage
\subsection{Use Case Description-22: Moderate Content \& Re-Index Ontology Taxonomy}
\label{subsec:uc22}

\begin{table}[H]
\centering
\small
\begin{tabularx}{\textwidth}{|l|X|}
\hline
\rowcolor{lightgray!40} \textbf{Use Case} & \textbf{UC22: Moderate Content \& Re-Index Ontology Taxonomy} \\ \hline
\textbf{Goal} & Administrator updates academic skill taxonomy, moderates custom tags, and triggers vector re-indexing. \\ \hline
\textbf{Preconditions} & Administrator is signed in with \texttt{ADMIN} superuser privileges. \\ \hline
\textbf{Success End Condition} & Taxonomy nodes updated; vector index refresh queued. Message: \textit{``Taxonomy updated and re-indexing initiated.''} \\ \hline
\textbf{Failed End Condition} & Operation failed. Message: \textit{``Unable to update taxonomy: Cyclic hierarchy detected.''} \\ \hline
\textbf{Primary Actors:} & System Administrator, Department Capstone Coordinator \\
\textbf{Secondary Actors:} & None \\ \hline
\textbf{Trigger} & Administrator opens the Taxonomy \& Content Moderation Console. \\ \hline
\textbf{Description / Main Success Scenario} &
\begin{enumerate}[leftmargin=*]
    \item Administrator opens the Academic Taxonomy Tree view.
    \item Administrator adds new technical skill entities (e.g., \texttt{LangChain}, \texttt{TinyML}) under parent domain nodes.
    \item Administrator reviews and approves or merges user-submitted custom skill tags.
    \item System verifies that the taxonomy remains an acyclic directed graph (DAG).
    \item Administrator clicks ``Re-Index Academic Embeddings''.
    \item System triggers background Celery worker to regenerate taxonomy vectors in Qdrant.
    \item System logs taxonomy modification event in administrative audit ledger.
\end{enumerate} \\ \hline
\textbf{Alternative Flows} &
\begin{itemize}[leftmargin=*]
    \item[\textbf{4.a.}] \textbf{Administrator attempts to create circular relationship:}
    \begin{itemize}[leftmargin=*]
        \item[4.a.1.] System rejects operation: \textit{``Circular dependency detected between domain nodes.''}
    \end{itemize}
\end{itemize} \\ \hline
\textbf{Quality Requirements} & 
1. Validate taxonomy DAG integrity within 100 milliseconds. \newline
2. Execute vector collection re-indexing asynchronously without interrupting search queries. \\ \hline
\end{tabularx}
\end{table}

\clearpage
\subsection{Use Case Description-23: Monitor Departmental Supervisory Allocation Dashboard}
\label{subsec:uc23}

\begin{table}[H]
\centering
\small
\begin{tabularx}{\textwidth}{|l|X|}
\hline
\rowcolor{lightgray!40} \textbf{Use Case} & \textbf{UC23: Monitor Departmental Supervisory Allocation Dashboard} \\ \hline
\textbf{Goal} & Coordinator monitors real-time faculty supervisory distribution, unserved student teams, and quota health. \\ \hline
\textbf{Preconditions} & Coordinator or auditor is signed in with administrative view permissions. \\ \hline
\textbf{Success End Condition} & Supervisory heatmap and unallocated student queues rendered. Message: \textit{``Departmental dashboard loaded.''} \\ \hline
\textbf{Failed End Condition} & Dashboard failed to load. Message: \textit{``Analytics data unavailable.''} \\ \hline
\textbf{Primary Actors:} & Department Capstone Coordinator, Institutional Auditor \\
\textbf{Secondary Actors:} & None \\ \hline
\textbf{Trigger} & Coordinator navigates to ``Departmental Allocation Console''. \\ \hline
\textbf{Description / Main Success Scenario} &
\begin{enumerate}[leftmargin=*]
    \item Coordinator opens the Department Allocation Dashboard.
    \item System aggregates real-time data across all active capstone teams and faculty rosters.
    \item System renders an interactive allocation heatmap showing active student load per faculty member.
    \item System highlights faculty members who are overloaded or have reached maximum quota.
    \item System displays a dedicated list of ``Unallocated Student Teams'' who have no confirmed advisor.
    \item Coordinator filters data by academic semester, program, or specialization track.
\end{enumerate} \\ \hline
\textbf{Alternative Flows} &
\begin{itemize}[leftmargin=*]
    \item[\textbf{5.a.}] \textbf{Coordinator identifies an unallocated team whose requests all expired:}
    \begin{itemize}[leftmargin=*]
        \item[5.a.1.] Coordinator uses administrative action to assign an available co-advisor directly.
    \end{itemize}
\end{itemize} \\ \hline
\textbf{Quality Requirements} & 
1. Load aggregate dashboard views for 200+ faculty and 500+ teams in under 1.5 seconds. \newline
2. Refresh data automatically every 60 seconds or on manual refresh click. \\ \hline
\end{tabularx}
\end{table}

\clearpage
\subsection{Use Case Description-24: Search and Export Institutional Compliance Audit Trail}
\label{subsec:uc24}

\begin{table}[H]
\centering
\small
\begin{tabularx}{\textwidth}{|l|X|}
\hline
\rowcolor{lightgray!40} \textbf{Use Case} & \textbf{UC24: Search and Export Institutional Compliance Audit Trail} \\ \hline
\textbf{Goal} & Trace system activities, review cryptographic transaction logs, and export verifiable audit records. \\ \hline
\textbf{Preconditions} & Administrator or institutional auditor has verified administrative credentials. \\ \hline
\textbf{Success End Condition} & Filtered audit records displayed; encrypted CSV/JSON audit file exported. Message: \textit{``Audit trail exported successfully.''} \\ \hline
\textbf{Failed End Condition} & Search or export failed. Message: \textit{``Audit search query failed: Invalid parameters.''} \\ \hline
\textbf{Primary Actors:} & System Administrator, Institutional Auditor, Department Coordinator \\
\textbf{Secondary Actors:} & None \\ \hline
\textbf{Trigger} & Auditor accesses the ``Audit Trail \& Governance'' portal. \\ \hline
\textbf{Description / Main Success Scenario} &
\begin{enumerate}[leftmargin=*]
    \item Auditor navigates to the Audit Trail interface.
    \item Auditor defines search criteria: Date Range, Actor Role, Event Type (e.g., \texttt{FSM\_STATE\_TRANSITION}), and User UUID.
    \item System queries the append-only audit database table.
    \item System validates cryptographic hash chaining across selected log records.
    \item System renders matching audit events with timestamp, user ID, IP address, and payload diff.
    \item Auditor clicks ``Export Signed Audit CSV''.
    \item System compiles audit dataset, attaches SHA-256 integrity signature, and serves download.
\end{enumerate} \\ \hline
\textbf{Alternative Flows} &
\begin{itemize}[leftmargin=*]
    \item[\textbf{4.a.}] \textbf{Cryptographic log chain mismatch detected (tampering attempt):}
    \begin{itemize}[leftmargin=*]
        \item[4.a.1.] System displays high-severity security alert and isolates affected record range.
        \item[4.a.2.] System dispatches immediate emergency security alert to System Administrators.
    \end{itemize}
\end{itemize} \\ \hline
\textbf{Quality Requirements} & 
1. Enforce append-only integrity; reject any UPDATE or DELETE operations on audit log tables. \newline
2. Export 10,000 audit log records to CSV in under 2 seconds. \\ \hline
\end{tabularx}
\end{table}

\clearpage
\subsection{Use Case Description-25: Generate Departmental Capstone Analytics \& Accreditation Report}
\label{subsec:uc25}

\begin{table}[H]
\centering
\small
\begin{tabularx}{\textwidth}{|l|X|}
\hline
\rowcolor{lightgray!40} \textbf{Use Case} & \textbf{UC25: Generate Departmental Capstone Analytics \& Accreditation Report} \\ \hline
\textbf{Goal} & Generate certified PDF and CSV institutional summary reports for accreditation boards (e.g., BAETE/IEB). \\ \hline
\textbf{Preconditions} & Authorized coordinator or auditor session active; semester capstone records synchronized. \\ \hline
\textbf{Success End Condition} & Certified institutional report generated and downloaded. Message: \textit{``Accreditation report ready for download.''} \\ \hline
\textbf{Failed End Condition} & Generation failed. Message: \textit{``Unable to compile report: Incomplete data for selected semester.''} \\ \hline
\textbf{Primary Actors:} & Department Capstone Coordinator, Institutional Auditor \\
\textbf{Secondary Actors:} & System Administrator \\ \hline
\textbf{Trigger} & Coordinator navigates to ``Reports \& Accreditation'' and clicks ``Generate Report''. \\ \hline
\textbf{Description / Main Success Scenario} &
\begin{enumerate}[leftmargin=*]
    \item Coordinator opens the Reports console.
    \item Coordinator selects report type: ``Semester Supervisory Allocation Summary'' or ``Interdisciplinary Research Velocity''.
    \item Coordinator selects academic term (e.g., ``Fall 2026'') and target department.
    \item Coordinator selects output format: PDF (formal executive report) or CSV (raw data).
    \item System aggregates metrics: total projects, supervisor-to-student ratios, multidisciplinary co-advisory percentage, and milestone adherence rates.
    \item System compiles the report document with institutional header, summary tables, and charts.
    \item System generates downloadable file and logs report generation in audit trail.
\end{enumerate} \\ \hline
\textbf{Alternative Flows} &
\begin{itemize}[leftmargin=*]
    \item[\textbf{5.a.}] \textbf{No finalized projects exist for selected term:}
    \begin{itemize}[leftmargin=*]
        \item[5.a.1.] System displays notice: \textit{``No completed records found for the selected semester.''}
    \end{itemize}
\end{itemize} \\ \hline
\textbf{Quality Requirements} & 
1. Render and compile multi-page PDF accreditation report within 3 seconds. \newline
2. Exclude private student contact details from aggregate institutional reports. \\ \hline
\end{tabularx}
\end{table}
